%% EXHALE: how the well-balanced reconstruction is tied to the hydrostatic
%% condition, and why it is what makes the LHS 1140 b base layer solvable.
%% Written 2026-09-13 from src/modules/states/Reconstruction.f90, the design
%% memo docs/well_balanced_flux_difference_design_20260910.md and the
%% measurements of docs/lhs1140b_stationary_L5c_20260913.md.
%% Author: Kwang-Il Seon
\documentclass[english,a4paper]{my_memo}
\synctex=1
\usepackage{booktabs}
\usepackage{amsmath}
\usepackage[authoryear,round]{natbib}
\usepackage{babel}
\emergencystretch=3em

\title{The well-balanced reconstruction and the hydrostatic condition}
\author{Kwang-Il Seon}
\date{Last updated: \docmoddate}

\begin{document}
\maketitle

\begin{abstract}
The \texttt{Well balanced:\ True} option of EXHALE is the reconstruction of
\citet{Kappeli2016}: every cell carries its own hydrostatic equilibrium, the
reconstruction and the Riemann solver act on the departure from it, and the
momentum source is the face-pressure difference of that same equilibrium, so
that at a discrete hydrostatic equilibrium the flux difference and the source
cancel in the algebra and not in floating point. This memo states the
connection term by term, maps it onto the code, and records why it is the one
change that lets the LHS\,1140\,b base layer (Mach $2.7\times10^{-7}$) reach a
certified stationary state: the ordinary reconstruction's truncation error in
the face pressure enters the HLLC contact speed as a spurious velocity one to
three decades above the physical one, and no solver can remove a residual the
scheme itself manufactures. Every number is measured on the tree of
2026-09-13 unless marked otherwise.
\end{abstract}

\section{Background: well-balanced schemes}

A scheme is called well balanced for a given steady state when its discrete
operator vanishes exactly on the discrete counterpart of that state, so that
the state is preserved to rounding and small departures from it are resolved
at the scheme's accuracy instead of being buried under the truncation error
of two large cancelling terms. The idea entered Godunov-type methods with
\citet{LeVeque1998}, who noted that a high-resolution method applied to a
balance law with a source term produces spurious waves near a quasi-steady
state unless the flux gradient and the source are discretized so as to cancel,
and proposed a quasi-steady wave-propagation algorithm that splits the cell
so that the source balances the flux difference inside it.
\citet{Botta2004} treated the case that matters here, nearly hydrostatic
flows of an atmosphere, where the pressure gradient and gravity are each many
orders larger than their difference: they showed that the standard
finite-volume discretization loses its accuracy in that regime and built a
well-balanced method around a local hydrostatic reconstruction of the pressure.

\citet{Kappeli2014} gave a well-balanced finite-volume scheme for the Euler
equations with a general gravitational potential by defining, inside each
cell, a local equilibrium at the cell's own specific entropy and integrating
the mechanical balance $h + \phi = {\rm const}$ (their eq.\ 2.7); density and
pressure are both extrapolated to the faces from that isentropic equilibrium
and the momentum source is the face-pressure difference of the same
equilibrium (their eqs.\ 2.14 and 2.26), which their theorem 1 shows to be a
second-order consistent discretization of $-\rho\,d\phi/dx$. Its limitation
is that the equilibrium preserved exactly is the isentropic one.
\citet{Kappeli2016} removed that limitation: the equilibrium is again the
mechanical balance, but the subcell integral is taken with a constant density
inside the cell, so no thermodynamic variable other than the cell's own
pressure and density enters, and ``the reconstruction preserves discrete
hydrostatic equilibria with arbitrary temperature or entropy stratification''
(their section 2.1.1). Density and velocity keep the standard reconstruction,
the pressure is reconstructed as a departure from the local equilibrium, and
the only requirement on the numerical flux is that it resolve a stationary
contact discontinuity exactly. That is the variant EXHALE implements
(decision 23 of 2026-09-10; \texttt{docs/well\_balanced\_flux\_difference\_design\_20260910.md}).
Provenance: the design memo of 2026-09-10 read the authors' ETH SAM
research reports (2015-40 and 2013-05, the submitted texts); the published
articles (\texttt{references/Kappeli\_2016A\&A\_587\_A94.pdf},
\texttt{references/Kappeli\_2014JCP\_259\_199.pdf}) were obtained on
2026-09-13 and every equation number quoted here and in the design memo was
checked against them: eqs.\ (10), (11), (13) to (18) and (24) to (27) of the
2016 article, eqs.\ (2.7), (2.9) to (2.14), (2.23), (2.25), (2.26) and
theorem 1 of the 2014 article carry the same numbers and the same content
as in the reports, and the sentence quoted from section 2.1.1 is verbatim.
The two papers were read in full; the other references of this section were
consulted for their scope and are cited for it.

Related constructions exist. \citet{Chandrashekar2015} give a second-order
well-balanced finite-volume scheme for the Euler equations with gravity that
preserves isothermal and polytropic hydrostatic states exactly with any
consistent numerical flux, by reconstructing the pressure and density as
departures from those particular equilibria. \citet{Edelmann2021} compare
well-balanced treatments of gravity in astrophysical (stellar) hydrodynamics,
where the same near-hydrostatic cancellation limits the accessible Mach
number, and show on stellar models that without such a treatment the spurious
velocities exceed the physical convective ones at low Mach number. The two
ingredients of the EXHALE scheme that the well-balanced reconstruction leans
on are older: the HLLC approximate Riemann solver of \citet{Toro1994}, which
restores the contact wave the HLL solver averages out and therefore resolves a
stationary contact exactly, and the WENO reconstruction in the efficient form
of \citet{Jiang1996}, whose third-order variant is the production
reconstruction of EXHALE (\texttt{Reconstruction scheme:\ WENO3}).

\section{The hydrostatic condition in a finite-volume scheme}

The stationary momentum equation of the wind is
\begin{equation}
  \frac{d}{dr}\!\left(\rho v^{2}\right) + \frac{dp}{dr}
  = -\rho\,\frac{d\phi}{dr} - \frac{2\rho v^{2}}{r},
  \label{eq:mom}
\end{equation}
and in the base layer of LHS\,1140\,b, where the flow is at Mach
$2.7\times10^{-7}$, it reduces to the hydrostatic balance
$dp/dr = -\rho\,d\phi/dr$ to one part in $10^{-13}$. The two terms of that
balance are each $4.4\times10^{5}$ times the wind's momentum flux in the same
layer (\emph{Update\_EXHALE} stage 1, section 143); the wind is their
difference.

A finite-volume scheme forms the two terms in different places. The pressure
enters through the numerical flux at the faces, $F_{j+1/2}$, built from face
states that a reconstruction (PLM or WENO3) extrapolates from the cell
averages; gravity enters as a cell-centered source, $-\rho_j\,\phi'(r_j)$. The
update of the momentum row is
\begin{equation}
  R_{2,j} = \frac{A_{j+1/2}F_{2,j+1/2} - A_{j-1/2}F_{2,j-1/2}}{V_j}
            + \rho_j\,\phi'(r_j),
\end{equation}
and at a hydrostatic column the two contributions cancel only up to the
truncation error of the reconstruction, of order $(\Delta r/H)^{p}$ times the
$O(1)$ terms, $p$ being the order of the reconstruction and $H$ the pressure
scale height. That residual is not small against the wind: on the LHS\,1140\,b
column the WENO3 face pressure differs from the cell pressure by
$5.7\times10^{-5}$ relative at $r = 1.001\,R_p$ (measured on the archived
state; $7.3\times10^{-6}$ after eight $1$-$2$-$1$ smoothings of the column
and $1.2\times10^{-7}$ on a ninth-order polynomial fit, converging at order
$2.97$ to $3.00$ under refinement), while the stationary velocity is
$2.7\times10^{-7}$ of the sound speed.

Where that error does its damage is the HLLC flux. Its contact speed is
\begin{equation}
  S^{*} = \frac{p_R - p_L + \rho_L v_L (S_L - v_L) - \rho_R v_R (S_R - v_R)}
               {\rho_L (S_L - v_L) - \rho_R (S_R - v_R)}
        \equiv S^{*}_{\rm vel} + S^{*}_{\rm pres},
  \qquad
  S^{*}_{\rm pres} = \frac{p_R - p_L}{\phi_L - \phi_R},
  \label{eq:sstar}
\end{equation}
and the face energy flux is $F_3 = S^{*}(E^{*} + p^{*})$. In the base cells
of LHS\,1140\,b, $S^{*}_{\rm vel}/c_s$ is $3.4\times10^{-7}$ to
$1.1\times10^{-6}$, the Mach number of the gas, while $S^{*}_{\rm pres}/c_s$ is
$5\times10^{-6}$ to $9\times10^{-4}$: the pressure-jump part is one to three
decades larger. The energy flux divergence of the base cells is then the
divergence of a spurious velocity, and the energy row reads $O(1)$ of its own
largest term in every cell below $1.05\,R_p$ (measured $-4.6\times10^{4}$
times the net heating at cell 2, $+2.7\times10^{3}$ at cell 5, $-52$ at cell
20, $+1.7$ at cell 41, where the velocity part alone is $+3.4$, $+1.1$, $+1.2$,
$+1.5$, i.e.\ within a factor two of the stationary balance). The seed does
not matter: three seeds whose base velocity differed by a factor $10^{4}$
gave the same energy maximum to every printed digit. A Newton solver cannot
remove a residual that the scheme manufactures from the pressure field alone.

\section{The well-balanced reconstruction}

\subsection{Each cell carries its own hydrostatic equilibrium}

Within cell $j$, hold the density at the cell value $\rho_j$ and integrate the
mechanical balance from the cell's own pressure:
\begin{equation}
  p_{{\rm eq},j}(r) = p_j - \rho_j\left[\phi(r) - \phi(r_j)\right].
  \label{eq:peq}
\end{equation}
This is eq.\ (16) of \citet{Kappeli2016}. It states nothing about the
temperature or the entropy, which is why it preserves an arbitrary
stratification, and it is second-order consistent with $dp/dr = -\rho\phi'$
in the cell width. Its two face values are the arrays \texttt{wb\_P\_up}(j)
and \texttt{wb\_P\_dn}(j) of \texttt{hydrostatic\_equilibrium\_of\_each\_cell}
(\texttt{src/modules/states/Reconstruction.f90}):
\begin{align}
  P_{\rm up}(j) &= p_j - \rho_j\left[\phi(r_{j+1/2}) - \phi(r_j)\right], \\
  P_{\rm dn}(j) &= p_j + \rho_j\left[\phi(r_j) - \phi(r_{j-1/2})\right].
\end{align}
The potential is the analytic function of $r$ the code already evaluates at
faces and centers (\texttt{Gphi\_i}, \texttt{Gphi\_c}), so the staggered
interpolation of the paper's eq.\ (15) is not needed.

\subsection{The reconstruction carries the departure}

The perturbation data of cell $j$ is the neighbor's pressure measured against
cell $j$'s equilibrium continued \emph{through the face} to the neighbor's
center, with cell $j$'s density up to the face and the neighbor's beyond it:
\begin{align}
  d_{+}(j) &= (p_{j+1} - p_j) + \rho_j\,\phi_{\rm up}(j) + \rho_{j+1}\,\phi_{\rm dn}(j+1), \\
  d_{-}(j) &= (p_{j-1} - p_j) - \rho_j\,\phi_{\rm dn}(j) - \rho_{j-1}\,\phi_{\rm up}(j-1),
\end{align}
with $\phi_{\rm up}(j) = \phi(r_{j+1/2}) - \phi(r_j)$ and
$\phi_{\rm dn}(j) = \phi(r_j) - \phi(r_{j-1/2})$, and the value at the cell's
own center identically zero. The limited slope (PLM) or the WENO3 stencil is
applied to the triple $(d_{-}, 0, d_{+})$ exactly as it is applied to
$(p_{j-1}, p_j, p_{j+1})$ otherwise, with the same limiter, the same smoothness
indicators and the same volume-share coefficients, and returns two face
departures. The face pressures handed to the Riemann solver are
$P_{\rm up}(j) + q_L(j)$ and $P_{\rm dn}(j+1) + q_R(j)$; density and velocity
keep their ordinary reconstruction.

The discrete hydrostatic equilibrium this preserves is eq.\ (18) of the paper,
\begin{equation}
  \frac{p_{j+1} - p_j}{\Delta r} = -\frac{\rho_j + \rho_{j+1}}{2}\,
  \frac{\phi_{j+1} - \phi_j}{\Delta r},
  \label{eq:discrete_eq}
\end{equation}
which is exactly the statement that the two neighboring extrapolations agree
at the shared face, $P_{\rm up}(j) = P_{\rm dn}(j+1)$: the face mismatch
\begin{equation}
  \delta p_{{\rm eq}}(j+1/2) = P_{\rm dn}(j+1) - P_{\rm up}(j) = d_{+}(j)
\end{equation}
(\texttt{wb\_dp\_eq} in \texttt{well\_balanced\_face\_departures}) vanishes on
it. On the discrete equilibrium the perturbation data is therefore
identically zero, the reconstruction returns the equilibrium itself, and the
face is a stationary contact discontinuity: equal pressure on both sides and
zero velocity. The only requirement on the flux is that it resolves one
exactly, $F([\rho_L, 0, p], [\rho_R, 0, p]) = [0, p, 0]$, which HLLC and Roe
do and LLF does not (the reason the key warns under \texttt{Numerical flux:\
LLF}). Continuing cell $j$'s own density over the whole gap to the neighboring
center instead of switching to the neighbor's beyond the face leaves data of
size $(\rho_j - \rho_{j+1})$ times the potential difference across half a
cell, which does not vanish on eq.\ \eqref{eq:discrete_eq}: measured, that
form leaves $2.5\times10^{-5}$ of the momentum weight on the discrete column
at $N = 250$ where the form above leaves $3.9\times10^{-14}$.

\subsection{The source is the same equilibrium}

The momentum source is no longer the cell-centered $-\rho_j\phi'(r_j)$ but
the face-pressure difference of the cell's own equilibrium,
\begin{equation}
  S_{2,j} = \frac{A_{j+1/2}P_{\rm up}(j) - A_{j-1/2}P_{\rm dn}(j)}{V_j}
  \quad\text{measured against the face pressures},
\end{equation}
which is a second-order consistent discretization of $-\rho\,d\phi/dr$
\citep[theorem 1]{Kappeli2014}; in the code \texttt{Source} returns
$S(2,j) = 0$ and the momentum row is assembled from the face pressures
measured against the equilibrium. At the discrete equilibrium the flux
difference and the source are then the same expression and cancel in the
algebra; away from it the departure carries the ordinary second-order
accuracy of the reconstruction, so the wind region is discretized as before.

\subsection{What is and is not preserved}

The equilibrium of eq.\ \eqref{eq:peq} holds the density constant inside a
cell. What the scheme preserves exactly is therefore the \emph{discrete}
equilibrium \eqref{eq:discrete_eq}, not the continuous hydrostatic solution
$p(r)$ with its own $\rho(r)$; the two agree to second order in the cell width.
The consequence is that the base layer of a well-balanced solution converges
with the grid at the same order as before, and the grid convergence of the
LHS\,1140\,b solution has not yet been measured (open in
\texttt{docs/PLAN\_20260913\_lhs\_stationary.md}, item L5c). Nothing thermal
is assumed: the temperature stratification of the base layer is whatever the
energy equation returns.

\section{What it did on LHS\,1140\,b}

The measurements are those of
\texttt{docs/lhs1140b\_stationary\_L5c\_20260913.md} on the reference case
(GJ\,1132 proxy spectrum, He/H $= 1.6261$, $K_{zz} = 10^{9}$\,cm$^2$\,s$^{-1}$,
scalar base at 226\,K and 1\,$\mu$bar), the archived 2026-08-30 solution
mapped onto the current grid as the seed, the partitioned stationary route
(\texttt{Restart intent:\ stationary}, \texttt{EXHALE\_PTC\_DTAU0=1.0}),
16 threads.

\begin{table}[h]
\centering
\caption{The hydrodynamic rows of outer pass 1 (cellwise maximum of
$|R|$ over the row's own largest term) and the outcome of the stationary
solve, same seed, input and binary.}
\begin{tabular}{lcccccl}
\toprule
option & hydro & mass & momentum & energy & s & outcome \\
\midrule
none & info$=2$ & $1.08$ & $3.2\times10^{-2}$ & $0.63$ & 1325 & pass refused \\
\texttt{Low-Mach damping:\ 2.0e-2} & info$=2$ & $1.32$ & $1.4\times10^{-2}$ & $1.24$ & 1232 & pass refused \\
\texttt{Well balanced:\ True} & info$=0$ & $2.6\times10^{-10}$ & $1.8\times10^{-14}$ & $4.6\times10^{-9}$ & 85 & CERTIFIED at pass 8 \\
\bottomrule
\end{tabular}
\end{table}

With the key on, the LHS\,1140\,b atomic wind with binary H/He element
diffusion is certified at outer pass 8 in about ten minutes (mass
$4.5\times10^{-10}$, momentum $2.3\times10^{-13}$, energy $5.1\times10^{-9}$,
gated element row $8.8\times10^{-6}$ of $10^{-5}$, face mass-flux spread
$1.5\times10^{-11}$). Three seeds whose base velocities differed by a factor
$10^{4}$ (the mapped archived state, the same with the stationary velocity
imposed below $1.27\,R_p$, and a Bernoulli-consistent temperature) reach one
solution: $\rho$ to $3.4\times10^{-10}$, $v$ to $2.5\times10^{-8}$, $\dot M$
to seven digits. The element relaxation, which refused every pass without the
key (\texttt{element\_step\_out\_of\_bounds} after 24 to 26 steps, the helium
fraction driven to unity by a hydrodynamic state whose mass flux fell 27
percent across the wind), reaches its fixed point in 32 steps in every pass
with it. \texttt{Low-Mach damping} does nothing to the energy row because it
is a fourth difference of the \emph{velocity} and does not touch the pressure
truncation error. On the HD\,209458\,b element-diffusion reload the key
changes nothing that matters: certified at the same outer pass 12, $\dot M$
moved by $1.3\times10^{-4}$, profiles beyond $2\,R_p$ by $2$ to
$3\times10^{-4}$, the base cells by 12 to 14 percent.

The solution itself differs from the archived one: $\log_{10}\dot M = 8.29$
against $7.80$, the base cell at 405\,K against 252\,K, wind densities three
to ten times higher. Part of that is the physics that moved since 2026-08-30
(a heating 3.4 times larger in the first cells after the photon-grid
quadrature correction, the Jupiter radius, the spectrum below 13.6\,eV), part
the base condition, and part the fact that the archived state was never
stationary below $1.1\,R_p$ under the present residual (the old gates began
at $1.2\,R_p$). Which part is which is the comparison of item L4 of the plan
and is not settled here.

\section{Why the key is off by default}

Turning it on changes the discretization of the pressure/gravity pair in
every run, not only in the base layer: the two discretizations are both
second-order consistent and neither is a subset of the other, so every golden
of the regression matrix would move (\texttt{docs/input\_schema.md} K46). The
decision of 2026-09-10 (decision 23 of
\texttt{docs/To\_be\_determined\_by\_user\_20260906.md}) made it an option,
stated in \texttt{EXHALE\_setup.out}, carried as the option token
\texttt{wellbal} in every state file, and left off. For LHS\,1140\,b it is
required: without it there is no certifiable state of the wind on the current
code, with any seed and by any route (plan items L1 to L5c). The decision of
2026-09-13 is to run every LHS\,1140\,b model of \texttt{LHS1140b/MODELS.md}
with the key on and to say so in each case's \texttt{REPRODUCE.md}. Whether
it should become the default for every planet is a separate question, to be
decided on the regression matrix and the Model A comparison, not here.

\bibliographystyle{aa}
\begin{thebibliography}{}
\bibitem[Botta et al.(2004)]{Botta2004}
Botta, N., Klein, R., Langenberg, S., \& L\"utzenkirchen, S.\ 2004, Journal of
Computational Physics, 196, 539, ``Well balanced finite volume methods for
nearly hydrostatic flows'' (ADS 2004JCoPh.196..539B)
\bibitem[Chandrashekar \& Klingenberg(2015)]{Chandrashekar2015}
Chandrashekar, P., \& Klingenberg, C.\ 2015, SIAM Journal on Scientific
Computing, 37, B382, ``A second order well-balanced finite volume scheme for
Euler equations with gravity'' (ADS 2015SJSC...37B.382C)
\bibitem[Edelmann et al.(2021)]{Edelmann2021}
Edelmann, P.~V.~F., Horst, L., Berberich, J.~P., Andrassy, R., Higl, J.,
Leidi, G., Klingenberg, C., \& R\"opke, F.~K.\ 2021, A\&A, 652, A53,
``Well-balanced treatment of gravity in astrophysical fluid dynamics
simulations at low Mach numbers'' (ADS 2021A\&A...652A..53E)
\bibitem[Jiang \& Shu(1996)]{Jiang1996}
Jiang, G.-S., \& Shu, C.-W.\ 1996, Journal of Computational Physics, 126,
202, ``Efficient implementation of weighted ENO schemes''
(ADS 1996JCoPh.126..202J)
\bibitem[K\"appeli \& Mishra(2014)]{Kappeli2014}
K\"appeli, R., \& Mishra, S.\ 2014, Journal of Computational Physics, 259,
199, ``Well-balanced schemes for the Euler equations with gravitation''
(ADS 2014JCoPh.259..199K; \texttt{references/Kappeli\_2014JCP\_259\_199.pdf},
and the submitted text ETH SAM Research Report 2013-05 beside it)
\bibitem[K\"appeli \& Mishra(2016)]{Kappeli2016}
K\"appeli, R., \& Mishra, S.\ 2016, A\&A, 587, A94, ``A well-balanced finite
volume scheme for the Euler equations with gravitation. The exact
preservation of hydrostatic equilibrium with arbitrary entropy
stratification'' (ADS 2016A\&A...587A..94K; \texttt{references/Kappeli\_2016A\&A\_587\_A94.pdf},
and the submitted text ETH SAM Research Report 2015-40 beside it)
\bibitem[LeVeque(1998)]{LeVeque1998}
LeVeque, R.~J.\ 1998, Journal of Computational Physics, 146, 346, ``Balancing
source terms and flux gradients in high-resolution Godunov methods: the
quasi-steady wave-propagation algorithm'' (ADS 1998JCoPh.146..346L)
\bibitem[Toro et al.(1994)]{Toro1994}
Toro, E.~F., Spruce, M., \& Speares, W.\ 1994, Shock Waves, 4, 25,
``Restoration of the contact surface in the HLL-Riemann solver''
(ADS 1994ShWav...4...25T)
\end{thebibliography}

\end{document}
